\part{Processes and scheduling}

\chapter{What a process is}

\begin{goals}
  \item The difference between a program and a process
  \item What context is and why it must be saved completely
  \item What information the process control block holds
\end{goals}

\section{Program versus process}

A \term{program} is a passive thing: bytes on a disk. A \term{process} is a program \emph{running},
plus everything that describes the state of that execution.

Run the same program twice and you get two processes. They share not a single byte of state: two
different stacks, two different heaps, two different sets of register values, two different
positions in the code. The program is a recipe; the process is a particular meal being cooked.

\begin{concept}{Context}
The \term{context} of a process is everything that must be saved and restored so that it can be
stopped and resumed without noticing. For a 32-bit x86 process that is:
\begin{itemize}
  \item the eight general-purpose registers,
  \item \reg{EIP}: where it was executing,
  \item \reg{EFLAGS}: its flag state,
  \item \reg{ESP} and \reg{SS}: its stack,
  \item \reg{CS} and \reg{DS}: its segments, which encode its privilege,
  \item \reg{CR3}: its address space.
\end{itemize}
Forget \emph{any} of these and the process resumes subtly wrong --- a comparison behaves
backwards, a pointer aims at nothing, a return goes to the wrong place. Context switching bugs are
notoriously hard to find precisely because the symptom appears long after the cause.
\end{concept}

\section{The process control block}

All of that lives in one structure, conventionally called the PCB:

\filetag{lytlnybl/include/tasks/procman.h}
\begin{lstlisting}[style=c]
typedef struct process {
    uint32_t pid_waiting_for;   /* pid of a child we are waiting on */
    uint32_t parent_pid;

    uintptr_t user_heap_end;    /* for sbrk */
    uint32_t heap_pages_allocated;

    uint32_t wake_tick;         /* for sleep */
    process_reading_state_t reading_state;  /* for a blocked read */

    uint32_t pid;
    process_registers_t regs;   /* THE CONTEXT */

    process_states_t state;
    process_type_t type;        /* kernel or user */

    struct process* next;       /* singly linked list */

    page_directory_t* page_directory;

    void* kstack;               /* kernel stack */
    void* ustack;               /* user stack */
} __attribute__((packed)) process_t;
\end{lstlisting}

\begin{warn}{\cod{packed} does not belong here}
\cod{\_\_attribute\_\_((packed))} is correct on structures that mirror a hardware layout --- the
TSS, GDT entries, filesystem records. \cod{process\_t} is purely internal: nothing outside the
kernel ever sees its byte layout.

Packing it forces unaligned field access, which is slower on x86 and would fault outright on ARM or
MIPS. It also makes taking a pointer to a field technically undefined behaviour. Removing
\cod{packed} here would be a small, safe improvement.
\end{warn}

\section{The five states}

\begin{center}
\begin{tikzpicture}[font=\sffamily\scriptsize,
  s/.style={draw=accent,fill=accentlight,rounded corners=3pt,minimum width=2.3cm,minimum height=0.9cm,align=center},
  a/.style={-{Stealth[length=2.2mm]},thick}]

  \node[s,fill=greenlight,draw=green] (run) at (0,0) {RUNNING};
  \node[s] (rdy) at (-4.4,0) {READY};
  \node[s,fill=amberlight,draw=amber] (blk) at (0,-2.4) {BLOCKED};
  \node[s,fill=purplelight,draw=purple] (slp) at (4.4,0) {SLEEPING};
  \node[s,fill=redlight,draw=red] (trm) at (0,2.4) {TERMINATED};

  \draw[a,green] (rdy) to[bend left=18] node[above,font=\tiny,text=gray]{scheduler picks it} (run);
  \draw[a,accent] (run) to[bend left=18] node[below,font=\tiny,text=gray]{time slice ends} (rdy);
  \draw[a,amber] (run) -- node[right,font=\tiny,text=gray]{read() with no data} (blk);
  \draw[a,amber] (blk) to[bend right=45] node[left,font=\tiny,text=gray]{key arrives} (rdy);
  \draw[a,purple] (run) -- node[above,font=\tiny,text=gray]{sleep(n)} (slp);
  \draw[a,purple] (slp) to[bend right=45] node[above,font=\tiny,text=gray]{timer expires} (rdy);
  \draw[a,red] (run) -- node[right,font=\tiny,text=gray]{exit() or fault} (trm);
\end{tikzpicture}
\end{center}

\begin{center}
\begin{tabularx}{\textwidth}{@{}l X@{}}
\toprule
\textbf{State} & \textbf{Meaning} \\
\midrule
\cod{RUNNING}    & Executing right now. Exactly one process is in this state \\
\cod{READY}      & Could run; waiting for its turn \\
\cod{BLOCKED}    & Waiting for an event: a keystroke, a child exiting \\
\cod{SLEEPING}   & Waiting for a specific time to arrive \\
\cod{TERMINATED} & Finished; its resources have not been reclaimed yet \\
\bottomrule
\end{tabularx}
\end{center}

The distinction between \cod{BLOCKED} and \cod{SLEEPING} is worth noting: a sleeping process is
woken by the timer, which knows exactly when; a blocked process is woken by whichever subsystem it
was waiting on, which does not know when. They are handled in different places --- sleeping in
\cod{timer\_handler}, blocked in \cod{keyboard\_handler} or in the \cod{exit} system call.

\section{Creating a process}

\filetag{lytlnybl/tasks/procman.c}
\begin{lstlisting}[style=c]
process_t* create_process(void* task_address, process_type_t type) {
    process_t* new_process = kmalloc(sizeof(process_t));

    new_process->wake_tick = 0;
    new_process->type = type;
    new_process->pid = next_pid++;
    new_process->state = PROCESS_READY;
    new_process->kstack = kmalloc(KERNEL_STACK_SIZE);   /* 16 KB */
    new_process->next = NULL;

    new_process->regs.eax = 0;  /* ... all registers zeroed ... */
    new_process->regs.eip = (uintptr_t)task_address;
    new_process->regs.eflags = 0x202;   /* interrupts enabled! */

    if (type == PROCESS_KERNEL) create_kprocess(new_process);
    else                        create_uprocess(new_process);

    /* append to the end of the list */
    process_t* traversal_process = process_head;
    while (traversal_process) {
        if (!traversal_process->next) {
            traversal_process->next = new_process;
            break;
        }
        traversal_process = traversal_process->next;
    }
    return new_process;
}
\end{lstlisting}

\begin{concept}{The process has not run yet, but its context already exists}
This is the elegant part. The new process is given a \emph{fabricated} context: registers at zero,
\reg{EIP} at its entry point, flags with interrupts on, a fresh stack pointer.

Nothing distinguishes that fabricated state from one produced by genuinely interrupting a running
process. So the scheduler needs no special ``start'' path: it simply restores this context like any
other, \cod{iret} runs, and the process begins. \hi{Starting a process and resuming one are
the same operation.}
\end{concept}

\subsection{Kernel processes versus user processes}

\begin{lstlisting}[style=c]
void create_kprocess(process_t* new_process) {
    new_process->regs.esp = (uint32_t)(new_process->kstack) + KERNEL_STACK_SIZE;
    new_process->regs.cs = 0x08;    /* kernel code, ring 0 */
    new_process->regs.ds = 0x10;    /* kernel data */
    new_process->regs.ss = 0x10;
    new_process->page_directory = kernel_directory;   /* shares the kernel's space */
}
\end{lstlisting}

A kernel process runs in ring 0, in the kernel's own address space, on a heap-allocated stack. It is
essentially a thread. Note \cod{esp} points at the \emph{end} of the allocation, because the stack
grows down.

A user process, by contrast, gets its own page directory, its own stack mapped at
\hexa{BFFFF000}, its own heap, and ring 3 selectors. Part X covers it.

% ============================================================
\begin{recipe}{creating kernel processes}{ch30-kernel-process}
Creates two kernel processes that do nothing but increment a counter, then prints both counters six
times from the main kernel process. All three advance, and none of them yields: the timer takes the
CPU away regardless. That is preemption, in about twenty lines.

The comments state the two rules for a kernel process: it must never return, because there is no
return address on its stack, and it cannot call \cod{sleep()} or \cod{exit()}, because those are
system calls and it is already inside the kernel.
\end{recipe}

% ============================================================
\chapter{Context switching}

\begin{goals}
  \item Understand how a process is swapped for another
  \item See why the interrupt stack frame is the key to the whole mechanism
  \item Know what the TSS has to do with it
\end{goals}

\section{The idea}

Recall from Chapter 21 that the interrupt stub passes C a pointer to the interrupted code's complete
saved state, and that \cod{iret} restores whatever that structure contains.

A context switch therefore needs no special machinery at all:

\begin{enumerate}
  \item Copy the current contents of \cod{registers\_t} into the current process's PCB.
  \item Copy the \emph{new} process's PCB into that same \cod{registers\_t}.
  \item Change \reg{CR3} to the new process's page directory.
  \item Return normally. \cod{iret} resumes the new process.
\end{enumerate}

\begin{center}
\begin{tikzpicture}[font=\sffamily\scriptsize,
  b/.style={draw=grayborder,rounded corners=3pt,minimum width=2.9cm,minimum height=2.0cm,align=center}]
  \node[b,fill=greenlight,draw=green] (a) at (0,0) {\textbf{PCB of A}\\{\tiny eip, esp, eax\ldots}};
  \node[b,fill=amberlight,draw=amber] (s) at (5,0) {\textbf{registers\_t}\\{\tiny on the kernel stack}};
  \node[b,fill=purplelight,draw=purple] (c) at (10,0) {\textbf{PCB of B}\\{\tiny eip, esp, eax\ldots}};

  \draw[-{Stealth[length=2.5mm]},green,thick] (s) to[bend right=25] node[below,font=\tiny]{1. save\_context} (a);
  \draw[-{Stealth[length=2.5mm]},purple,thick] (c) to[bend right=25] node[below,font=\tiny]{2. load\_context} (s);
  \draw[-{Stealth[length=2.5mm]},red,thick] (5,-1.6) -- (5,-2.6) node[below,font=\tiny,text=red]{3. iret runs B};
\end{tikzpicture}
\end{center}

\section{The code}

\filetag{lytlnybl/tasks/context.c}
\begin{lstlisting}[style=c]
void context_switch(process_t* old_process, process_t* new_process,
                    registers_t* regs) {
    current_process = new_process;

    if (old_process->state == PROCESS_RUNNING) {
        old_process->state = PROCESS_READY;
    }

    save_context(old_process, regs);
    new_process->state = PROCESS_RUNNING;

    return_to_user = (new_process->type == PROCESS_USER);

    if (new_process->type == PROCESS_USER) {
        tss.esp0 = (uintptr_t)(new_process->kstack) + KERNEL_STACK_SIZE;
        set_cr3((uintptr_t)new_process->page_directory);
    } else if (new_process->type == PROCESS_KERNEL) {
        set_cr3((uintptr_t)kernel_directory);
    }

    if (old_process->state == PROCESS_TERMINATED) {
        destroy_process(old_process);
    }

    load_context(new_process, regs);
}
\end{lstlisting}

\subsection{Why the state is checked before overwriting it}

The first \cod{if} only demotes \cod{RUNNING} to \cod{READY}. That matters: if the process arrived
here because it called \cod{exit} (state \cod{TERMINATED}), or because it blocked on a read (state
\cod{BLOCKED}), or because it slept, that state must be preserved. Unconditionally setting
\cod{READY} would make a blocked process immediately runnable again, and the \cod{read} would return
garbage.

\subsection{Saving the stack correctly}

\filetag{lytlnybl/tasks/context.c}
\begin{lstlisting}[style=c]
void save_context(process_t* process, registers_t* regs) {
    process->regs.eip = regs->eip;
    process->regs.cs  = regs->cs;
    process->regs.eflags = regs->eflags;
    process->regs.ds  = regs->ds;
    /* ... the eight general-purpose registers ... */

    if (process->type == PROCESS_KERNEL) {
        process->regs.esp = regs->esp;        /* from pusha */
    } else {
        process->regs.esp = regs->user_esp;   /* pushed by the processor */
        process->regs.ss  = regs->ss;
    }
}
\end{lstlisting}

This is the conditional from Chapter 21 made concrete. For a kernel process there was no ring change,
so the only \reg{ESP} available is the one \cod{pusha} saved. For a user process the processor
pushed the real user \reg{ESP} and \reg{SS} onto the kernel stack, and those are the ones that
matter --- the \cod{pusha} value is the kernel stack pointer, which is meaningless to the user
program.

\begin{danger}{Getting this backwards}
Save \cod{regs->esp} for a user process and on resumption it would run with a \emph{kernel} stack
pointer while in ring 3. The first push would fault --- or worse, if that address happened to be
mapped user-accessible, it would silently corrupt kernel memory. This is exactly the class of bug
that appears ``randomly'', hours later, in an unrelated subsystem.
\end{danger}

\subsection{\cod{tss.esp0}: preparing for the next interrupt}

\begin{lstlisting}[style=c]
tss.esp0 = (uintptr_t)(new_process->kstack) + KERNEL_STACK_SIZE;
\end{lstlisting}

When a ring 3 process is interrupted, the processor reads \cod{esp0} from the TSS and switches to
that stack automatically. So before running a user process, the kernel must point \cod{esp0} at
\emph{that process's} kernel stack.

Forget this line and every process would share one kernel stack: a nested interrupt would clobber
another process's saved state. Part X returns to the TSS in detail.

\section{The scheduler}

\filetag{lytlnybl/tasks/scheduler.c}
\begin{lstlisting}[style=c]
void schedule(registers_t* regs) {
    scheduler_tick_count++;
    if (scheduler_tick_count >= time_slice) {
        scheduler_tick_count = 0;
        process_t* next_process = get_next_process();
        if (next_process != current_process) {
            context_switch(current_process, next_process, regs);
        }
    }
}
\end{lstlisting}

\cod{schedule} is called on every timer tick, but only switches every \cod{time\_slice} ticks. That
constant is the \term{quantum}: how long a process gets before being preempted.

\begin{concept}{The quantum is a trade-off}
Short quantum: the system feels responsive, but more time is spent switching than working --- every
switch costs a save, a restore, a \reg{CR3} reload and a full TLB flush.

Long quantum: less overhead, but a process hogging the CPU makes everything else feel frozen.

Real schedulers vary it dynamically and add priorities: interactive processes get short, frequent
slices; compute-bound ones get long, rare ones. lytlnyblOS uses a fixed quantum, which is the right
choice for something you want to be able to read.
\end{concept}

\subsection{Round robin}

\begin{lstlisting}[style=c]
process_t* get_next_process() {
    process_t* traversal_process = current_process;
    do {
        if (traversal_process->state == PROCESS_READY) {
            return traversal_process;
        }
        if (!traversal_process->next) {
            traversal_process = process_head;   /* wrap around */
        } else {
            traversal_process = traversal_process->next;
        }
    } while (!(traversal_process->state == PROCESS_RUNNING));

    return current_process;
}
\end{lstlisting}

Start at the current process, walk forward, wrap at the end, return the first \cod{READY} one.
Everyone gets a turn, in order, with no favourites. It is the simplest fair policy that exists.

\begin{danger}{What happens when nobody is ready}
The loop stops when it comes back round to a \cod{RUNNING} process, and then returns
\cod{current\_process}. But consider: the shell calls \cod{read} and blocks. There is no other
runnable process. \cod{get\_next\_process} walks the whole list, finds nothing \cod{READY}, comes
back to the shell --- which is now \cod{BLOCKED}, not \cod{RUNNING} --- and the loop condition
\cod{while (!(state == RUNNING))} keeps it spinning.

In practice the kernel's own initial process is always \cod{RUNNING}, so the loop terminates and
returns it. But that is luck, not design: the exit condition depends on a process that happens to
sit in an infinite loop at the end of \cod{kernel\_main}.

The standard fix is an \hi{idle process}: a kernel process, always \cod{READY}, that does
nothing but \cod{hlt} in a loop. It guarantees the scheduler always has somewhere to go, and it
lets the CPU actually sleep when there is no work --- which the current \cod{for(;;)} at the end of
\cod{kernel\_main} does not.
\end{danger}

\section{Destroying a process}

\filetag{lytlnybl/tasks/procman.c}
\begin{lstlisting}[style=c]
void destroy_process(process_t* proc) {
    if (!proc || proc->state == PROCESS_RUNNING) return;

    /* unlink it from the list */
    process_t* traversal_process = process_head;
    while (traversal_process) {
        if (traversal_process->pid == process_head->pid
            && !traversal_process->next
            && traversal_process->pid == proc->pid) {
            process_head = traversal_process->next;
            break;
        }
        if (!traversal_process->next) { traversal_process = NULL; break; }
        if (traversal_process->next->pid == proc->pid) {
            traversal_process->next = traversal_process->next->next;
            break;
        }
        traversal_process = traversal_process->next;
    }
    if (!traversal_process) return;

    kfree(proc->kstack);
    if (proc->type == PROCESS_USER) {
        unmap_page(proc->page_directory, USER_STACK_TOP - 4096);
        unmap_page(proc->page_directory, USER_CODE_BASE);
        unmap_page(proc->page_directory, USER_VGA);
        unmap_page(kernel_directory, (uintptr_t)(proc->page_directory));
    }
    kfree(proc);
}
\end{lstlisting}

This function has several real problems, and it is worth being explicit about them because they are
representative of how cleanup code goes wrong.

\begin{danger}{Three bugs in one function}
\textbf{1. The head is never removed if it has a successor.} The first condition requires
\cod{!traversal\_process->next} --- that the head is also the \emph{last} node. If the head has a
successor, that branch fails, the second returns early, and the third only ever inspects
\cod{next}. A terminated head process stays in the list forever.

\textbf{2. \cod{unmap\_page(kernel\_directory, ...)} damages every process.} Every process copies the
kernel's directory entries, which means they all point at the \emph{same} kernel page tables.
Removing a mapping from \cod{kernel\_directory} removes it from the shared table --- so it
disappears from every address space at once, permanently punching a hole in the identity map.

\textbf{3. No physical memory is ever reclaimed.} \cod{unmap\_page} clears the present bit but never
calls \cod{free\_frame}. The page directory, the code frames, the stack frame and the heap frames
all stay marked used forever. Run enough programs and the machine runs out of RAM with nothing
running.
\end{danger}

\begin{concept}{And a fourth, subtler one}
\cod{destroy\_process} is called from inside \cod{context\_switch}, which is running \emph{on the
kernel stack of the process being destroyed}. \cod{kfree(proc->kstack)} marks that stack free while
the code is still using it, and it continues using it until the \cod{iret}.

Nothing breaks today because nothing allocates in between. But it is a textbook use-after-free, and
the correct pattern is deferred cleanup: mark the process as reapable and free its resources later,
from a context that is not standing on them.
\end{concept}

\begin{tryit}[Watch the teardown problem]
Boot the system and run \cod{run /bin/user\_test} five or six times in a row from the shell. The
program works each time, but the display progressively degrades --- output from different runs
interleaves and overwrites earlier lines. That is the cleanup path leaving the system in a
progressively worse state. A correct teardown would leave the system indistinguishable from before.
\end{tryit}

\begin{summary}
The kernel can now create, schedule, switch and (imperfectly) destroy processes. Everything so far
runs in ring 0, though. The next part does the thing this whole course has been building towards:
dropping a program into ring 3 and giving it a controlled way back in.
\end{summary}
